\chapter{Environnement et données}

\section{Configuration de référence}

\begin{table}[H]
\centering
\caption{Composants de l'environnement de test}
\begin{tabularx}{\textwidth}{|L{3.4cm}|L{4.1cm}|Y|}
\hline
\rowcolor{ctmint}\textbf{Composant} & \textbf{Version de référence} & \textbf{Rôle} \\ \hline
Système & Windows, fuseau Africa/Tunis & Poste de développement \\ \hline
PHP & 8.5.8 & Exécution Laravel et PHPUnit \\ \hline
Laravel & 13.19.0 & API REST et logique métier \\ \hline
PostgreSQL & 18.0, 64 bits & Base locale de développement \\ \hline
Node.js & 22.19.0 & Frontend et outils simulateur \\ \hline
pnpm & 11.16.0 & Gestion des dépendances frontend \\ \hline
React / Vite & React 19.2.7, Vite 8.2.1 & Interface web et build \\ \hline
Python & 3.13 dans l'image de test & Gateway OCPP et pytest \\ \hline
Docker & 29.6.1 & Simulateurs OCPP et paiement \\ \hline
Redis & 7 Alpine & Cache, sessions, files et diffusion \\ \hline
\end{tabularx}
\end{table}

\section{Services}

\begin{table}[H]
\centering
\caption{Services attendus pour une recette verticale}
\begin{tabularx}{\textwidth}{|L{3.2cm}|L{4.3cm}|Y|}
\hline
\rowcolor{ctmint}\textbf{Service} & \textbf{Adresse locale} & \textbf{Commande} \\ \hline
Pile Docker complète & Ports configurés dans \path{infra/.env} & \texttt{pnpm stack:up} \\ \hline
Frontend & \texttt{http://localhost:5173} & \texttt{pnpm dev:frontend} \\ \hline
Backend REST & \texttt{http://localhost:8000} & \texttt{pnpm dev:backend} \\ \hline
Worker & Processus sans port & \texttt{pnpm dev:queue} \\ \hline
Scheduler & Processus sans port & \texttt{pnpm dev:scheduler} \\ \hline
Reverb & \texttt{ws://localhost:8080} & \texttt{pnpm dev:reverb} \\ \hline
Gateway OCPP & \texttt{ws://localhost:9000} & \texttt{pnpm ocpp:up} \\ \hline
Paiement simulé & \texttt{http://localhost:9090} & \texttt{pnpm payment:up} \\ \hline
Documentation API & \url{http://localhost:8000/docs/api} & Incluse dans le backend \\ \hline
Production & \url{https://chargetrackr.me} & Déploiement Azure, TLS géré par Caddy \\ \hline
\end{tabularx}
\end{table}

\section{Profils de test}

Les comptes doivent être créés par seed ou par les workflows d'invitation. Les
adresses suivantes sont des exemples ; leurs mots de passe ne doivent jamais
figurer dans le dépôt ou dans une preuve.

\begin{table}[H]
\centering
\caption{Jeux d'acteurs synthétiques}
\begin{tabularx}{\textwidth}{|L{3.2cm}|L{4.2cm}|Y|}
\hline
\rowcolor{ctmint}\textbf{Rôle} & \textbf{Périmètre} & \textbf{Utilisation} \\ \hline
Super Administrateur & Plateforme entière & Demandes de démo, organisations, offres SaaS, intégrations et audit. \\ \hline
Administrateur A & Organisation Tunis Network Ops & Employés, clients, tarifs, facturation et rapports de l'organisation A. \\ \hline
Opérateur A & Organisation A & Bornes, connecteurs, alertes, affectations et maintenance. \\ \hline
Technicien A & Organisation A & Interventions et maintenances assignées. \\ \hline
Administrateur B & Organisation Sahel Charge Network & Vérification d'isolation par rapport à l'organisation A. \\ \hline
Client 1 & Compte global & Sessions et abonnements auprès de plusieurs organisations. \\ \hline
Visiteur & Non authentifié & Inscription, connexion et demande de démonstration. \\ \hline
\end{tabularx}
\end{table}

\section{Précautions}

\begin{itemize}[leftmargin=6mm]
  \item Réinitialiser les simulateurs avant une campagne reproductible.
  \item Utiliser des références de borne uniques.
  \item Ne jamais tester Resend avec une liste réelle de destinataires.
  \item Conserver les captures, réponses HTTP et journaux sans cookies, clés,
        secrets OAuth, mots de passe ou identifiants de paiement.
  \item Noter le commit, la date, le navigateur et les services actifs pour
        chaque campagne manuelle.
\end{itemize}
